\begin{textsample}{NEG Dim 2 – ce – Score 1.00 – ce/ce\_inf\_12.txt}  \label{ex:f2_neg_210}
ALGUNS PONTOS QUE NÃO PODEMOS PERDER DE VISTA NA NOSSA PRÁTICA PEDAGÓGICA AO DISCUTIRMOS O QUADRO SÍNTESE DA INTEGRAÇÃO E ORGANIZAÇÃO DAS EXPERIÊNCIAS GARANTIR OS DIREITOS DE APRENDIZAGEM E DESENVOLVIMENTO ( Brincar, Conviver, Participar, Explorar, Expressar e Conhecer -se ) Os direitos de aprendizagem e desenvolvimento precisam ser garantidos e devem ser concretizados nas experiências previstas nas DCNEI/2009 e na BNCC/2017.
Não podem ser considerados de forma fragmentada e ganham especificidades nos diferentes campos de experiência.
CAMPOS DE EXPERIÊNCIAS ( O eu, o outro e o nós ; Corpo, gestos e movimentos ; Traços, sons, cores e formas ; Escuta, fala, pensamento e imaginação ; Espaços, tempos, quantidades, relações e transformações. ) Constituem um “ arranjo curricular ” que partem das experiências das crianças, de suas ações cotidianas e abrigam seus saberes e os conhecimentos, entrelaçando aos conhecimentos que fazem parte ao patrimônio cultural.
EXPERIÊNCIAS Têm as interações e a brincadeira como eixos norteadores, previstas nas DCNEI/2009 ( Incisos Art. 9º ) e nos Campos de Experiência - BNCC/2017 ( Campos de Experiência com seus Objetivos de Aprendizagem e Desenvolvimento ).
APRENDIZAGENS POSSÍVEIS Ao participarem de experiências significativas, em que seus direitos de aprendizagem e desenvolvimento são garantidos, as crianças aprendem e se desenvolvem.
PONTO DE PARTIDA PARA A ORGANIZAÇÃO DA AÇÃO PEDAGÓGICA Interesses e especificidades das crianças, identificados a partir da observação e registro de suas ações.
ORIENTAÇÕES DIDÁTICO-METODOLÓGICAS QUE CONSIDEREM POSSIBILIDADES DE : a ) situações de interação ( criança/crianças ; professora/professora/criança e crianças ) ; b ) variedade de brincadeiras e desafios ; c ) escolhas e produção pelas crianças ; d ) escuta e respeito aos seus interesses e ritmos ; e ) relação dialógica e negociada ; f ) ação criativa, exploratória e representativa das crianças em diversas linguagens.
ORGANIZAÇÃO Das crianças, de acordo com seus próprios arranjos, da rotina, do tempo, espaço e materiais.
FAIXA ETÁRIA - Bebês : 0 a 1 e 6 meses - Crianças bem pequenas : 1 e 7 meses a 3 anos e 11 meses - Crianças pequenas : 4 anos a 5 anos e 11 meses DIREITOS DE APRENDIZAGEM E DESENVOLVIMENTO Os direitos de aprendizagem e desenvolvimento precisam ser garantidos e devem ser concretizados nos campos de experiência.
Não podem ser considerados de forma fragmentada e ganham especificidades nos diferentes campos de experiência.
CAMPOS DE EXPERIÊNCIAS Constituem um arranjo curricular que acolhe as situações e as experiências concretas da vida cotidiana das crianças e seus saberes, entrelaçando aos conhecimentos que fazem parte ao patrimônio cultural.
OBJETIVOS DE APRENDIZAGEM DESENVOLVIMENTO As aprendizagens essenciais compreendem tanto comportamentos, habilidades e conhecimentos quanto vivências que promovem aprendizagens e desenvolvimento nos diversos campos de experiências, sempre tomando as interações e a brincadeira como eixos norteadores.
ORGANIZAÇÃO E INTEGRAÇÃO DAS EXPERIÊNCIAS Referem -se à organização de práticas pedagógicas elaboradas com base na escuta da criança, respeitando as culturas infantis e as demais práticas culturais e considerando os princípios da didática do fazer : ludicidade, continuidade e significatividade ( BONDIOLI ; MANTOVANI, 1998 ).
A organização e integração das experiências incluem as orientações metodológicas que preveem diversificadas possibilidades de interação ( criança/crianças ; professora/professora ( e outros/outras profissionais da instituição)/criança e crianças entre si ) ; de escolhas e produção pelas crianças ; de escuta e respeito aos seus interesses e ritmos ; de diálogo e negociação ; diversidade de brincadeiras, situações desafiadoras envolvendo formas diferentes de representação ( em diversas linguagens ) que incentivem a ação criativa e exploratória das crianças.

% matched lemmas: 
\end{textsample}
